\begin{ledgroup}
\footnotesize
\pstart
\noindent
Die nachfolgenden Stücke N.\;\ref{41667} und N.\;\ref{55606}
dokumentieren Leibnizens erste Auseinandersetzung 
mit den Thesen des Jesuiten
\indexN{\textls{Vanni}, Giovanni Francesco 1638\textendash1709}Giovanni Francesco Vannis
über ein Grundprinzip der Statik,
die eine über die Grenzen Italiens hinaus heftig ausgetragene Kontroverse entfachten
(siehe dazu die Monographie \cite{02126}\textsc{M.~Torrini}, \textit{Dopo Galileo}, Florenz 1979)
und mit denen sich Leibniz über die folgende Jahrzehnte mehrmals befassen sollte (siehe neben N.\;\ref{intro_vanni-specimen} die Stücke N.\;\ref{intro_vanni-demonstratio}, \ref{41199}, \ref{intro_vanni-solutio-problematis}, und \ref{intro_vanni-investigatio} und \ref{intro_theoremata-duo-vanni}).
\pend
\pstart
\indexN{\textls{Vanni}, Giovanni Francesco 1638\textendash1709}Vanni veröffentlichte um 1684 in Rom die anonyme Flugschrift \cite{01995}\textit{Specimen libri de momentis gravium} als Auszug und Vorankündigung seines späteren Buchs \cite{02000}\textit{Exegeses physico-mathematicae} (Rom 1685).
Im \textit{Specimen} griff Vanni den bekannten, von Stevin, Galilei und Torricelli prominent verfochtenen 
Satz der Zerlegung von Kräften auf schiefen Ebenen an. 
Seine kontroverse These und der polemische Ton zogen die Aufmerksamkeit der italienischen und europäischen Gelehrtenwelt schlagartig auf sich.
\pend
\pstart
Dank der Vermittlung 
seines Förderers
\indexN{\textls{Magliabechi} (Magliabecchi, Magliabecki, Magliabekius, Magliabequius), Antonio 1633\textendash1714}A.~Magliabechi
und \indexN{\textls{Papebroch} (Papenbrockius), Daniel 1628\textendash1714}Daniel Papebrochs
gelang es Vanni, das \cite{01995}\textit{Specimen} in den \cite{01023}\textit{Acta Eruditorum} von November 1684 veröffentlichen zu lassen.
Der Mitherausgeber der \cite{01023}\textit{Acta},
\indexN{\textls{Pfautz} (Pfauzius), Christoph 1645\textendash1711}C.~Pfautz,
schickte allerdings zunächst den Text an Leibniz (\cite{02006}Brief vom 17./27. September 1684, \textit{LSB} III,~4 N.\;70), von dem er sich eine inhaltliche Einordnung und ggf. eine kritische Stellungnahme, die den Aufsatz begleiten sollte, erbat.
Leibniz verfasste eine solche im Konzept N.\;\ref{41667}; dessen Entstehung ist sicher zwischen dem 27. September 1684 und, wie aus folgenden Umständen ersichtlich wird, dem 17. Oktober 1684 zu verorten.
Denn an diesem Tag schickte Leibniz dem anderen Mitherausgeber der \textit{Acta}, \indexN{\textls{Mencke} (Menken, Menkenius, Menque), Otto 1644\textendash1707}Otto Mencke,
eine gekürzte Fassung von N.\;\ref{41667}, in der er seine Kritik an \indexN{\textls{Vanni}, Giovanni Francesco 1638\textendash1709}Vanni viel zurückhaltender als im Konzept äußerte: \glqq breviter et mollissimis verbis\grqq\ (\cite{02007}\textit{LSB} I,~4 N.\;398).
Diese Fassung der Bemerkung (von der keine handschriftlichen Textzeugen erhalten sind) 
erschien anonym in den \textit{Acta Eruditorum} von November 1684 (S.\;514)
im Anschluss an das \cite{01995}\textit{Specimen}.%
\pend
\pstart
\indexN{\textls{Vanni}, Giovanni Francesco 1638\textendash1709}Vannis Reaktion auf die Kritik ist uns durch seine Briefe an \indexN{\textls{Magliabechi} (Magliabecchi, Magliabecki, Magliabekius, Magliabequius), Antonio 1633\textendash1714}Magliabechi bekannt und wurde in dieser Form um 1685 auch an Leibniz herangetragen (siehe N.\;\ref{32310}). Sie veranlasste eine zweite Phase seiner Auseinandersetzung mit Vannis Thesen, der im Komplex um die \textit{Demonstratio geometrica} von 1685 ihren Niederschlag fand (N.\;\ref{intro_vanni-demonstratio}).
Des Weiteren verfasste Leibniz zwischen Ende Oktober 1686 und Mitte Februar 1687 auf Wunsch der Herausgeber der \textit{Acta Eruditorum} eine Rezension von Vannis angekündigtem Buch \textit{Exegeses}, nachdem es um 1685 erschienen war (N.\;\ref{41199}).
\pend
\end{ledgroup}




